\section{Reproducibility Details}
\label{app:reproducibility}

The main paper reports model identifiers and reasoning-effort settings; this
appendix records the operational path used for the reported sweeps. The full
step-by-step workflow is maintained in \path{docs/REPRODUCIBILITY.md}, and each
run also writes its exact configuration to
\path{data/runs/<run_id>/run_config.json}.

\paragraph{OpenAI Codex invocation.}
\textsc{GPT-5.5} runs used the OpenAI Codex CLI, authenticated through the
author's ChatGPT subscription. The adapter invokes this command shape:
\begin{lstlisting}
codex exec --skip-git-repo-check --ephemeral -s read-only \
  -m gpt-5.5 -i <image_path> -o <last_message_path> \
  -C <temporary_workdir> --color never \
  -c reasoning_effort=<medium|extra_high>
\end{lstlisting}
The image is attached once per request, and \texttt{-o} captures the final
message for normalization. The paper sweeps used two retries with exponential
backoff and per-sample timeouts of 1800 seconds for \texttt{medium} and 2400
seconds for \texttt{extra\_high}.

\paragraph{Claude invocation.}
\textsc{Claude Opus 4.7} (1M context) runs used the Claude Code CLI,
authenticated through the author's Claude subscription. The adapter invokes this
command shape:
\begin{lstlisting}
claude --print --add-dir <image_parent> \
  --model claude-opus-4-7[1m] --effort <high|max> \
  --output-format text --no-session-persistence
\end{lstlisting}
The target image is referenced in the prompt as \texttt{@<image\_path>}, and
\texttt{--add-dir} grants read access to the sample directory. The paper sweeps
used two retries with exponential backoff and per-sample timeouts of 1800
seconds for \texttt{high} and 2400 seconds for \texttt{max}.

\paragraph{Artifacts.}
Both multimodal paths receive the same zero-shot prompt and are normalized by the
same prediction normalizer before parsing and rendering. Run artifacts include
the raw response, normalized response, latency, adapter metadata, full evaluation
result, and aggregate summaries. The non-LLM baselines use the same runner and
artifact schema but do not call an external model-serving interface.
